
This diagnostic study characterizes two mechanistically distinct pathways by which WGA improvement occurs on Waterbirds WILDS. DFR achieves WGA = 0.91 without changing the backbone (layer4 cosine similarity = 1.000000 $\pm$ 1e-14 with ERM, confirmed at all 3 seed pairs). GroupDRO achieves WGA = 0.88 through a three-step mechanistic chain: minority group upweighting (5.01\% of training data, exponentiated gradient ascent) $\rightarrow$ group-balanced gradient propagating through all backbone layers (backbone/head L2 ratio = 6.47--7.03) $\rightarrow$ significantly reduced background linear decodability in layer4 features (ERM 0.9838 $\rightarrow$ GroupDRO 0.9530, p = 0.0039, Cohen's d = 6.48, CONFIRMED).

The backbone-vs-head typology matters for method selection and future method design. The probe protocol established here --- frozen layer4 features, sklearn L-BFGS C=1e9, full test set, paired t-test across matched seeds --- is a low-cost, reproducible backbone spurious encoding diagnostic.

Open questions remain: the suppression-vs-dilution mechanistic ambiguity requires directional subspace analysis; the WGA correlation needs per-seed WGA measurements for full confirmation; layer-wise spurious encoding profiles would characterize where GroupDRO's backbone effect is localized; and core attribute probe accuracy would confirm that useful feature encoding is preserved. Two pathways to WGA improvement exist in the Waterbirds WILDS setting; understanding when and why each operates is a prerequisite for principled method selection.

